\subsection{A row-independent P5 action across a cut}

\begin{proposition}[Proved: the relative sign obstruction]
\label{CP2:A-P5-row-independent}
Let $S$ be a finite nonempty left-$N$, right-$P_5$ biset
with no right fixed state. Suppose the two right band-row
actors agree on every state:
\[
 se_{0j}=se_{1j}\qquad(s\in S,\ j=0,1).
\]
Let $A$ be the finite left $N$ set of regular column-core
orbits. If the ordinary left $N$ act cone on $A$ has zero
rational degree-two products, then a collage on $S$ whose
second rational cohomology has dimension at most one has
zero degree-two squares.

The hypothesis on that act cone holds for every $N$ of
order at most seven, or for arbitrary finite $N$ when
$|A|\le2$. No bounded-homology hypothesis is required.
\end{proposition}
\begin{proof}
Let $U\subseteq S$ be the full band core, $X_U\subseteq X$
the subcollage, and $E_U\subseteq E$ their two-sided
action bars. Since there is no fixed state, every orbit
in $U$ is a regular two-point column orbit.
The full-core nullhomotopy gives
\[
 X_U\simeq BP_5\vee
 \operatorname{hocofib}\bigl(B(N,A)\longrightarrow BN\bigr).
\]
Thus every rational degree-two product on $X_U$ is zero
under the stated hypothesis.

The row interchange $\sigma(e_{ij})=e_{1-i,j}$, fixing
the other corner and every state, defines an involution
$F$ of the entire pair $(X,X_U)$. The row-independence
hypothesis is precisely the compatibility needed here.
We compute its relative action, retaining the actual
subcollage rather than replacing it by an abstract complex.

Tensor the length-two left projective $P_5$ resolution
from Lemma~\ref{CP2:A-order-five-corner} with the right
module $T=\mathbb Q[S]/\mathbb Q[U]$.
Every band element annihilates $T$. The degree-zero and
degree-two terms therefore vanish, and the degree-one
term is $T$. The semilinear lift used in
Lemma~\ref{CP2:A-mixed-five-reductions} acts by $-1$
on this surviving term. The relative right-action bar
is consequently $T[1]$, with $F=-1$, as an $N$ complex.
This computes the relative bar because each projective
resolution term is flat and the original permutation-module
sequence $0\to\mathbb Q[U]\to\mathbb Q[S]\to T\to0$
is exact. Taking derived $N$ coinvariants preserves this
relative comparison and the sign.

The two mapping-cylinder models share both corner spaces;
their quotient is $\Sigma(E/E_U)$. Hence $F=-1$ on
every rational relative homology group $H_*(X,X_U)$.
Over $\mathbb Q$ the eigenspaces of an involution split
exactly, so the homology sequence of the pair shows that
\[
 H_4(X_U;\mathbb Q)\longrightarrow H_4(X;\mathbb Q)
\]
surjects onto the $+1$ eigenspace. This conclusion does
not require $BN$ to be rationally acyclic.

If $H^2(X;\mathbb Q)$ has dimension one, a generator is
an eigenvector of $F^*$ with eigenvalue $+1$ or $-1$.
Its square has eigenvalue $+1$. It pairs to zero with
the negative fourth-homology eigenspace by its sign, and
with the positive eigenspace because that space comes
from $X_U$, where all second products vanish.
Its square is therefore zero. If the second cohomology
is zero the assertion is immediate.

Finally, all ordinary act cones over $|N|\le7$ have zero
second products by Corollary~\ref{CP2:A-high-seven-act}.
For $|A|\le2$ use the unconditional binary act theorem
(and the singleton cone). This proves both stated
unconditional instances.
\end{proof}

\begin{corollary}[Proved: scope refinement for mixed-five eight-state actions]
\label{CP2:A-P5-mixed-eight-row-scope}
In Proposition~\ref{CP2:A-P5-mixed-eight-scope}, right-act
indices $1,4,10,20,23$ may also be omitted. The remaining
twenty-four right acts against the forty-two mixed-five
targets give exactly $1008$ finite scopes.
\end{corollary}
\begin{proof}
Literal inspection of the checked normal-form maps shows
that the indicated five acts have equal row actors, and
the opposite corner has order five. The proposition
therefore applies. The other right indices retain their
previous scope. The new reduced enumeration must record
this gate explicitly; it does not alter or promote any
earlier uncompleted producer stream.
\end{proof}
